\chapter{Liquidation Waterfalls, Vesting, and Founder Contracts}
\label{ch:founder-contracts}
\chapterline

Chapter~\ref{ch:cap-tables-safes} built the share-count algebra of cap tables across seed SAFEs and priced venture rounds. Yet two companies with identical cap tables -- say Founders at $55\%$ and Investors or a Corporate Parent at $45\%$ -- can produce completely different cash outcomes when the company is sold, when a co-founder leaves in Month~4, or when an engineer takes a below-market salary to build a spin-out subsidiary for an established business.

A share count only measures your pro-rata weight inside one class of stock. Corporate charters, employment agreements, and shareholders' agreements wrap those shares in priority waterfalls, vesting forfeiture schedules, tax elections, salary--equity trade-offs, and related-party covenants. In this chapter we analyze the four contract arenas every technical founder or early partner faces: venture preferred liquidation waterfalls, co-founder vesting and tax elections, founder cash salary versus equity from first principles, and minority contracts when building a venture with an existing operating business.

\section{Preferred Stock and Liquidation Waterfalls}
\label{sec:liquidation-waterfalls}

Suppose an investor wires \$5{,}000{,}000 into a startup for $20\%$ of the company at a \$25{,}000{,}000 post-money valuation, while the two founders hold the remaining $80\%$. Two weeks later, before spending a dollar of the new capital, the founders accept an acquisition offer for \$10{,}000{,}000 in cash (\$5{,}000{,}000 for the code plus the \$5{,}000{,}000 sitting in the corporate bank account).

If the investor held ordinary Common Stock alongside the founders, an $80 / 20$ split of the \$10{,}000{,}000 proceeds would hand the founders \$8{,}000{,}000 and return only \$2{,}000{,}000 to the investor -- allowing the founders to pocket \$3{,}000{,}000 of the cash the investor wired fourteen days earlier. To prevent that moral hazard, cash venture investors purchase \textbf{Convertible Preferred Stock}, which sits ahead of Common Stock in the \textbf{liquidation waterfall} (the priority rule governing how net equity proceeds $E \ge 0$, after paying debt creditors and transaction expenses, are distributed upon an acquisition, merger, or winding up).

\subsection{1x Non-Participating Preferred and the Conversion Threshold}
\label{subsec:non-participating-preferred}

The clean market baseline in venture financing is \textbf{1x Non-Participating Preferred Stock}. Suppose a single investor class invests cash $I > 0$ with a $1\times$ preference multiple, creating a \textbf{liquidation preference} $L = 1 \cdot I = I$, and holds fraction $\alpha \in (0, 1)$ of fully diluted common-equivalent shares.

At exit value $E \ge 0$, ``non-participating'' means the investor chooses \emph{either} to take their preference \emph{or} to convert to common, whichever pays more:
\begin{enumerate}
  \item \textbf{Remain Preferred}: collect up to $L$ ahead of Common Stock, and nothing further; or
  \item \textbf{Convert to Common Stock}: waive the preference $L$ completely and collect pro-rata fraction $\alpha E$ alongside Common Stock.
\end{enumerate}

The investor's payout $P_{\text{inv}}^{\text{NP}}(E)$ and the residual payout $P_{\text{com}}^{\text{NP}}(E) = E - P_{\text{inv}}^{\text{NP}}(E)$ to Common Stock (founders and option holders) follow a three-regime piecewise linear function:
\begin{equation}\label{eq:non-participating-waterfall}
  P_{\text{inv}}^{\text{NP}}(E) = \min\!\bigl(E,\; \max(L, \alpha E)\bigr) =
  \begin{cases}
    E, & 0 \le E \le L, \\
    L, & L < E \le \dfrac{L}{\alpha}, \\[0.55em]
    \alpha E, & E > \dfrac{L}{\alpha},
  \end{cases}
  \qquad
  P_{\text{com}}^{\text{NP}}(E) =
  \begin{cases}
    0, & 0 \le E \le L, \\
    E - L, & L < E \le \dfrac{L}{\alpha}, \\[0.55em]
    (1 - \alpha)E, & E > \dfrac{L}{\alpha}.
  \end{cases}
\end{equation}

Setting $\alpha E^* = L$ locates the exact \textbf{conversion threshold}:
\begin{equation}\label{eq:conversion-threshold}
  E^* = \frac{L}{\alpha}.
\end{equation}
For exits $E \le L$, the investor recovers every available dollar while common receives zero. Across the catch-up zone $L < E \le L/\alpha$, the investor's payout stays flat at $L$ while $100\%$ of every marginal dollar goes to common shareholders. Above $E^* = L/\alpha$, preferred converts to common and every shareholder is paid pro rata. When $L = I$ and no subsequent round has diluted $\alpha = I / V_{\text{post}}$, we have $E^* = I / (I / V_{\text{post}}) = V_{\text{post}}$: the investor converts in every sale above their post-money entry valuation.

\subsection{Participating Preferred and Capped Participation}
\label{subsec:participating-preferred}

A more aggressive structure replaces ``or'' with ``and.'' Under \textbf{Uncapped Participating Preferred} (the \textbf{double dip}), the investor first takes $L$ off the top \emph{and then also} takes pro-rata share $\alpha$ of the remaining proceeds $\max(0, E - L)$:
\begin{equation}\label{eq:participating-waterfall}
  P_{\text{inv}}^{\text{P}}(E) = \min\!\bigl(E,\; L + \alpha \max(0, E - L)\bigr) =
  \begin{cases}
    E, & 0 \le E \le L, \\
    L + \alpha(E - L), & E > L,
  \end{cases}
\end{equation}
leaving common holders $P_{\text{com}}^{\text{P}}(E) = (1 - \alpha)\max(0, E - L)$. Rewriting the upper branch of~\eqref{eq:participating-waterfall} as
\[
  L + \alpha(E - L) = \alpha E + (1 - \alpha)L
\]
shows that at every exit $E > L$, uncapped participating preferred transfers a fixed surcharge $(1 - \alpha)L$ directly from common holders to the preferred investor, who never voluntarily converts in an acquisition (converting only if forced by an automatic-conversion threshold in a qualified IPO).

Under \textbf{Capped Participating Preferred} with return multiple cap $k > 1$ (for example, a $3\times$ cap limiting the preference-plus-participation track to $kI$), the investor collects $L + \alpha(E - L)$ until reaching $kI$ at $E_{\text{cap}} = L + \frac{kI - L}{\alpha}$, stays flat at $kI$ across $E_{\text{cap}} < E \le \frac{kI}{\alpha}$, and converts fully to common to collect $\alpha E$ once $E > \frac{kI}{\alpha}$.

\subsection{Multiple Liquidation Preferences and Stacked Seniority}
\label{subsec:multiple-preferences-seniority}

In late-stage or turbulent markets, two terms frequently turn high paper valuations into traps:
\begin{itemize}
  \item \textbf{Multiple Liquidation Preferences ($m\times$, with $m > 1$)}: setting $L = mI$ (such as $2\times$ or $3\times$ cash invested) requires returning two or three times the investor's check before junior tiers collect a dollar.
  \item \textbf{Stacked Seniority vs.\ Pari Passu}: across multiple rounds, \textbf{standard seniority} (last-in, first-out) pays Series~C in full before Series~B, Series~B before Series~A, and Series~A before Common. Under \textbf{pari passu} priority, preferred classes share a single senior tier and divide insufficient proceeds in proportion to their preference entitlements $L_j / \sum_r L_r$.
\end{itemize}

\begin{example}[How a \$200M Headline Valuation Wipes Founders to \$0 in an \$80M Acquisition]
Suppose a startup raises three rounds of venture capital:
\begin{itemize}
  \item \textbf{Series~A}: \$5M invested at \$25M post-money ($1\times$ non-participating, $L_A = \text{USD } 5\text{M}$).
  \item \textbf{Series~B}: \$15M invested at \$75M post-money ($1\times$ non-participating, senior to A, $L_B = \text{USD } 15\text{M}$).
  \item \textbf{Series~C}: Growth slows, yet the founders chase a ``\$200M unicorn-style valuation'' headline. An investor agrees to invest \$30M for $\alpha_C = 15.0\%$ on condition of a \textbf{$2\times$ senior non-participating preference} ($L_C = 2 \times \text{USD } 30\text{M} = \text{USD } 60\text{M}$, senior to Series~B and Series~A).
\end{itemize}
Total cash ever raised is \$50M, and post-Series-C fully diluted ownership stands at Series~C $= 15.0\%$, Series~B $= 17.0\%$, Series~A $= 13.6\%$, and Founders plus Early Employees (Common) $= 54.4\%$. Eighteen months later, an acquirer buys the company for \textbf{\$80{,}000{,}000} in cash -- \$30M above total capital invested. Trace the senior waterfall from top to bottom:
\begin{enumerate}
  \item \textbf{Tier 1 -- Series~C ($2\times$ Senior)}: converting to common would pay Series~C at most $0.15 \times 80 = 12$ million dollars (\$12M), far below $L_C = \text{USD } 60\text{M}$. Series~C stays Preferred and collects \textbf{\$60M}. Remaining pool: \$20M.
  \item \textbf{Tier 2 -- Series~B ($1\times$ Senior to A)}: converting into the remaining \$20M pool pays far less than $L_B = \text{USD } 15\text{M}$. Series~B stays Preferred and collects \textbf{\$15M}. Remaining pool: \$5M.
  \item \textbf{Tier 3 -- Series~A ($1\times$ Senior to Common)}: Series~A is owed $L_A = \text{USD } 5\text{M}$ and collects the final \textbf{\$5M}. Remaining pool: \textbf{\$0}.
  \item \textbf{Tier 4 -- Common Stock (Founders and Early Employees, $54.4\%$ ownership)}: \textbf{\$0}.
\end{enumerate}
Had the founders accepted a clean \$100M post-money valuation at $1\times$ preference ($L_C = \text{USD } 30\text{M}$, total preferences \$50M), an \$80M sale would have left roughly \$30M for common and converting holders. Chasing the \$200M headline with a $2\times$ senior preference raised the preference stack to \$80M and wiped $54.4\%$ of the cap table down to zero. Figure~\ref{fig:waterfall-payout-curves} plots investor and common payouts across all three preference structures.
\end{example}

\begin{figure}[htbp]
\centering
\begin{tikzpicture}
  % Left panel: Investor Payout P_inv(E) for $20M invested for alpha = 25% ($80M post)
  \begin{axis}[
    name=invplot,
    width=6.0cm, height=5.1cm,
    xmin=0, xmax=120,
    ymin=0, ymax=65,
    xlabel={Exit Proceeds $E$ (USD M)},
    ylabel={Investor Payout $P_{\text{inv}}(E)$ (USD M)},
    title={Investor Payout ($I=\text{USD }20\text{M},\;\alpha=25\%$)},
    title style={font=\footnotesize\bfseries},
    label style={font=\scriptsize},
    tick label style={font=\scriptsize},
    grid=major,
    grid style={gray!22},
    legend style={font=\tiny, at={(0.03,0.97)}, anchor=north west, draw=gray!40, fill=white, inner sep=2pt},
  ]
    \addplot[sourceblue, very thick] coordinates {(0,0) (20,20) (80,20) (120,30)};
    \addlegendentry{1x Non-Part.}
    \addplot[quotientorange, very thick, dashed] coordinates {(0,0) (20,20) (120,45)};
    \addlegendentry{1x Participating}
    \addplot[trapred, very thick, dotted] coordinates {(0,0) (40,40) (120,40)};
    \addlegendentry{2x Non-Part.}
    \fill[sourceblue] (axis cs:80,20) circle (1.6pt);
    \node[font=\tiny, text=sourceblue, anchor=north west] at (axis cs:76,19) {$E^*\!=\!80$};
  \end{axis}

  % Right panel: Common / Founder Payout P_com(E) = E - P_inv(E)
  \begin{axis}[
    at={(invplot.east)}, anchor=west, xshift=1.15cm,
    width=6.0cm, height=5.1cm,
    xmin=0, xmax=120,
    ymin=0, ymax=95,
    xlabel={Exit Proceeds $E$ (USD M)},
    ylabel={Common Payout $P_{\text{com}}(E)$ (USD M)},
    title={Common / Founder Payout ($1-\alpha=75\%$)},
    title style={font=\footnotesize\bfseries},
    label style={font=\scriptsize},
    tick label style={font=\scriptsize},
    grid=major,
    grid style={gray!22},
    legend style={font=\tiny, at={(0.03,0.97)}, anchor=north west, draw=gray!40, fill=white, inner sep=2pt},
  ]
    \addplot[sourceblue, very thick] coordinates {(0,0) (20,0) (80,60) (120,90)};
    \addlegendentry{1x Non-Part.}
    \addplot[quotientorange, very thick, dashed] coordinates {(0,0) (20,0) (120,75)};
    \addlegendentry{1x Participating}
    \addplot[trapred, very thick, dotted] coordinates {(0,0) (40,0) (120,80)};
    \addlegendentry{2x Non-Part.}
  \end{axis}
\end{tikzpicture}
\caption{Exit waterfall curves for a \$20M investment buying $\alpha = 25\%$ (\$80M post-money). Under $1\times$ non-participating preferred, common catches up dollar-for-dollar across $E \in [20, 80]$ and collects $75\%$ above $E^* = \text{USD } 80\text{M}$; uncapped participation deducts $(1-\alpha)L = \text{USD } 15\text{M}$ from common at every higher exit, while a $2\times$ preference zeroes common out up to \$40M.}
\label{fig:waterfall-payout-curves}
\end{figure}

\subsection{Anti-Dilution Protection in Down Rounds}
\label{subsec:anti-dilution-down-rounds}

Preferred Stock carries an original issue price $p_{\text{orig}}$ and a \textbf{conversion price} $CP$ (initially $CP_1 = p_{\text{orig}}$), converting into $p_{\text{orig}} / CP$ shares of Common Stock per preferred share. If the company later sells shares in a \textbf{down round} at $p_{\text{down}} < CP_1$, price-based anti-dilution reduces $CP_1$ to $CP_2 < CP_1$, lifting each preferred share's common-conversion ratio above $1$ without requiring additional cash from the prior investor.

The market-standard formula is \textbf{Broad-Based Weighted Average Anti-Dilution}:
\begin{equation}\label{eq:broad-based-weighted-average}
  CP_2 = CP_1 \cdot \frac{A + B}{A + C},
\end{equation}
where $A$ is the pre-round fully diluted share count (including common, as-converted preferred, and the full reserved option pool), $B = I_{\text{down}} / CP_1$ is the number of shares the down-round cash $I_{\text{down}}$ \emph{would} have bought at the old conversion price $CP_1$, and $C = I_{\text{down}} / p_{\text{down}}$ is the actual number of shares issued in the down round ($C > B$). Because $A$ anchors both numerator and denominator, a small bridge round produces a proportionally small adjustment: with $A = 10{,}000{,}000$, $CP_1 = \text{USD } 2.00$, $I_{\text{down}} = \text{USD } 1{,}000{,}000$, and $p_{\text{down}} = \text{USD } 1.00$, we have $B = 500{,}000$ and $C = 1{,}000{,}000$, giving $CP_2 = 2.00 \times \frac{10.5}{11.0} = \text{USD } 1.9091$.

\begin{trapbox}
\textbf{Predatory Full-Ratchet Anti-Dilution.}
A \textbf{Full Ratchet} clause ignores round size and resets $CP_2 = p_{\text{down}}$ unconditionally, even if the company sells only \$100{,}000 of shares at $p_{\text{down}}$. In the example above, if Series~A originally invested \$5{,}000{,}000 at \$2.00 for $2{,}500{,}000$ shares, raising a \$1{,}000{,}000 bridge at \$1.00 under Broad-Based Weighted Average grants Series~A $119{,}048$ bonus common-equivalent shares; under a Full Ratchet, $CP_2$ drops straight to \$1.00, doubling Series~A from $2{,}500{,}000$ to $5{,}000{,}000$ common-equivalent shares ($+2{,}500{,}000$ free shares) at the direct expense of founders and employees. Always require broad-based weighted-average anti-dilution, never full ratchet.
\end{trapbox}

\section{Co-Founder Equity: Vesting, Section 83(b), IP, and Acceleration}
\label{sec:cofounder-vesting}

\subsection{Why Every Founder Must Vest}
\label{subsec:why-founders-vest}

Suppose two engineers incorporate a startup, issue $4{,}000{,}000$ unconditional shares to each ($50\%$ / $50\%$), and skip vesting because they trust each other. In Month~4, Co-Founder~B burns out, accepts a salaried big-tech job, and stops working on the startup -- walking out the door owning $50.00\%$ of the voting stock forever.

Co-Founder~A is now trapped carrying $50\%$ \textbf{dead equity}: working years to build a business for half the remaining upside while institutional investors refuse to wire seed or Series~A money until the departed co-founder surrenders most of that idle block. Mutual founder vesting protects the remaining builders when a co-founder partnership ends early: shares are earned month by month through ongoing work rather than handed over permanently at incorporation.

\subsection{Four-Year Reverse Vesting With a One-Year Cliff}
\label{subsec:reverse-vesting-cliff}

The standard schedule is \textbf{four-year vesting with a one-year cliff} across $T = 48$ months on a founder grant $S_F$:
\begin{equation}\label{eq:vesting-schedule}
  S_{\text{vested}}(t) =
  \begin{cases}
    0, & 0 \le t < 12, \\[0.35em]
    \dfrac{\floor{t}}{48}\,S_F, & 12 \le t < 48, \\[0.45em]
    S_F, & t \ge 48,
  \end{cases}
\end{equation}
where $t$ is elapsed months of service. Leaving before Month~12 yields zero vested shares; completing Month~12 vests $\frac{12}{48} = 25.00\%$ at once, followed by $\frac{1}{48}$ ($2.0833\%$) per completed month through Month~48.

Rather than receiving options, founders purchase actual shares on Day~1 under a \textbf{Restricted Stock Purchase Agreement (RSPA)} with \textbf{reverse vesting}: the founder holds and votes all $S_F = 4{,}000{,}000$ shares immediately, subject to a lapsing company \textbf{repurchase option} on the unvested portion $S_F - S_{\text{vested}}(t)$ at original par cost (\$0.00001 per share) upon termination of service.

\subsection{The Section 83(b) Election and Global Tax Equivalents}
\label{subsec:section-83b-election}

Across major tax systems (Section~83 of the US Internal Revenue Code, UK Section~431 ITEPA restricted-securities elections, and similar statutes globally), default tax law waits until restricted shares \emph{vest} and taxes the spread $(\text{FMV}_t - p_{\text{paid}})$ on every vesting date as \textbf{ordinary employment cash income}, even though private startup shares cannot be sold to pay the tax.

Under US Section~83(b) (and the 14-day UK Section~431 equivalent), a founder has an unextendable statutory deadline of \textbf{30 calendar days from the share purchase date} to elect taxation at \emph{grant} ($t = 0$) rather than at vesting. Compare the two outcomes on $S_F = 4{,}000{,}000$ unvested shares bought on Day~1 at $p_0 = \text{USD } 0.00001$ (\$40.00 total cash) when $\text{FMV}_0 = \text{USD } 0.00001$:

\begin{itemize}
  \item \textbf{Timely Section~83(b) filed within 30 calendar days}: Day~1 spread is $4{,}000{,}000 \times (0.00001 - 0.00001) = 0$ dollars, so immediate tax is \textbf{\$0}. Across Months~12--48 as Seed and Series~A rounds lift 409A Common FMV to \$0.25 or \$1.50 per share, tax at every monthly vest is \textbf{\$0}. Years later upon an actual cash acquisition (say at \$5.00/share = \$20{,}000{,}000), the entire gain is taxed once -- when cash arrives -- at long-term capital-gains rates (or up to $0\%$ federal tax if eligible under US Section~1202 Qualified Small Business Stock).
  \item \textbf{Forgot to file Section~83(b) within 30 days}: suppose Seed and Series~A financings push the average 409A Common FMV across the $4{,}000{,}000$ vesting shares during Years~1--4 to \$0.25001 per share. Across those monthly vests, the tax authority records $4{,}000{,}000 \times 0.25 = 1{,}000{,}000$ dollars of ordinary compensation income. At a $40\%$ combined federal, state, and payroll marginal rate, the founder owes \textbf{\$400{,}000 in personal cash taxes} on illiquid paper shares (while the startup owes employer payroll taxes) -- potentially bankrupting the founder years before any exit, with no ordinary-income tax refund if the company later fails.
\end{itemize}

\subsection{Single-Trigger vs.\ Double-Trigger Acceleration}
\label{subsec:acceleration-on-acquisition}

If the startup is acquired in Month~24 when half of your shares remain unvested, two acceleration structures exist:
\begin{itemize}
  \item \textbf{Single-Trigger Acceleration}: $100\%$ of unvested shares vest automatically upon closing of the acquisition. Acquirers resist single-trigger founder acceleration because software acquisitions rely on retaining core technical leadership post-close; if founders vest $100\%$ at closing and can walk out the next day, buyers carve a new retention pool out of the purchase price paid to existing shareholders.
  \item \textbf{Double-Trigger Acceleration} (the balanced market standard): unvested equity continues vesting on schedule after closing, and accelerates $100\%$ if \emph{both} (\emph{Trigger~1}) an acquisition closes \emph{and} (\emph{Trigger~2}) within 12 to 18 months the founder is terminated \textbf{without Cause} or resigns for \textbf{Good Reason} (constructive dismissal via material pay cut, forced relocation, or stripping of engineering authority).
\end{itemize}

\subsection{IP Assignment and Employer Contamination}
\label{subsec:ip-assignment}

Every founder and engineer must sign a \textbf{Proprietary Information and Inventions Agreement (PIIA)} containing \emph{present-tense assignment} language (``hereby assigns,'' rather than ``agrees to assign'') vesting code and patent title cleanly in the corporation, together with a \textbf{Prior Inventions Schedule} listing pre-existing personal libraries or projects before work starts. Never write a line of startup code, run a training job, or store architecture notes on a current employer's laptop, corporate VPN, cloud account, or paid working hours; verify your employment contract and governing state law (such as California Labor Code Section~2870) before writing Day~1 code so prior-employer contamination never clouds title during diligence.

\section{Founder Compensation: Salary vs.\ Equity From First Principles}
\label{sec:founder-salary-vs-equity}

Should founders draw a cash salary, work unpaid, or take a below-market wage? Engineers routinely get confused here because folklore mixes together three completely different situations: two penniless co-founders bootstrapping pre-seed, a funded Series~A startup with \$6{,}000{,}000 in the bank, and an operating subsidiary owned $85\%$ by a profitable corporate parent. First-principles accounting resolves all three with one conservation model.

\subsection{Sweat Capital and Forgone Market Compensation}
\label{subsec:sweat-capital}

Suppose an engineer can command annual cash compensation $w_{\text{mkt}} > 0$ in the open labor market (say \$220{,}000 per year), and instead works full-time on a new venture drawing actual annual cash pay $w_{\text{act}} \in [0, w_{\text{mkt}}]$. Every year, that founder contributes an implicit economic input -- \textbf{sweat capital} -- equal to their \textbf{forgone market compensation}:
\begin{equation}\label{eq:forgone-salary}
  \Delta w = w_{\text{mkt}} - w_{\text{act}} \ge 0.
\end{equation}
Economically, working for $w_{\text{act}}$ is identical to being paid your full market wage $w_{\text{mkt}}$ in cash and immediately wiring $\Delta w$ back into the company's treasury every payroll period alongside any direct cash capital $K \ge 0$ you invested out of pocket.

\subsection{When Founders Are Unpaid vs.\ Paid Across Company Stages}
\label{subsec:unpaid-vs-paid-stages}

Tracking $\Delta w$ alongside corporate cash balance $C_0$ explains why rational founder pay shifts across stages:

\begin{enumerate}
  \item \textbf{Pre-Funding / Bootstrapped Stage ($C_0 \approx 0$)}: before customer gross profit or outside capital arrives, paying salaries is physically impossible without personal cash injections. If Co-Founder~A and Co-Founder~B have similar market opportunity costs ($w_{\text{mkt}} \approx \text{USD } 200{,}000$), both work full-time at $w_{\text{act}} = 0$, and both split incorporation costs equally ($K_A = K_B$), their ongoing economic inputs ($\Delta w_A = \Delta w_B$) are symmetric, justifying a $50 / 50$ equity split (subject to 4-year vesting). Conversely, if Co-Founder~A works full-time unpaid ($\Delta w_A = \text{USD } 200{,}000$/yr) and puts in $K_A = \text{USD } 60{,}000$ of personal savings while Co-Founder~B keeps a full-time salaried job working weekends ($\Delta w_B \approx \text{USD } 35{,}000$/yr, $K_B = 0$), a $50 / 50$ split violates conservation of contribution.
  \item \textbf{Funded Seed and Series~A+ Stage ($C_0 > 0$)}: once a startup raises institutional Seed or Series~A capital (or generates sufficient gross profit), \emph{disciplined venture investors actively insist that founders draw a livable baseline salary} -- typically enough to cover rent, healthcare, and family living costs comfortably without personal debt, while remaining below big-tech peak executive cash packages. Why? Because an unpaid founder watching personal savings hit zero faces acute personal solvency pressure: if a competitor dangles a premature \$8{,}000{,}000 acquisition that puts \$2{,}500{,}000 cash into a desperate founder's bank account, the founder is rationally tempted to sell early rather than compound the company toward \$250{,}000{,}000. On the financial statements (Part~I), founder salary plus employer payroll taxes flows straight through Income Statement Operating Expenses (R\&D / G\&A) and increases monthly net burn $B(t)$ in equation~\eqref{eq:net-burn}: paying two founders \$145{,}000 each (roughly \$325{,}000/year loaded) consumes \$487{,}500 over an 18-month seed runway, which is why seed founder salaries are calibrated to personal sustainability first and step up toward market executive medians at Series~A and Series~B.
\end{enumerate}

\subsection{The Risk-Adjusted Salary--Equity Trade-Off}
\label{subsec:risk-adjusted-salary-equity}

Why is \$100{,}000 of forgone cash salary $\Delta w$ \emph{not} fairly compensated by \$100{,}000 of nominal startup paper equity? Because market salary is guaranteed, immediate, liquid cash, whereas early-stage common equity is locked up for $T$ years and succeeds in producing a liquid common payout with probability $p_s \in (0, 1]$ (failing or getting wiped out by preferences with probability $1 - p_s$).

Let $r > 0$ be the annual opportunity discount rate over an exit horizon of $T$ years. Defining the \textbf{risk-and-illiquidity multiplier}
\begin{equation}\label{eq:risk-multiplier}
  \rho = \frac{(1 + r)^T}{p_s} > 1,
\end{equation}
sacrificing annual cash $\Delta w$ over $\tau$ years (cumulative forgone cash $\tau\,\Delta w$) has non-negative expected present value if and only if the founder's common ownership fraction $\omega$ and conditional exit value $E_{\text{com}}$ (net to common upon success) satisfy the \textbf{participation inequality}:
\begin{equation}\label{eq:salary-equity-tradeoff}
  \frac{p_s\,\omega\,E_{\text{com}}}{(1 + r)^T} \ge \tau\,\Delta w
  \quad\Longleftrightarrow\quad
  \omega\,E_{\text{com}} \ge \rho\,\tau\,\Delta w,
  \qquad
  E_{\text{BE}} = \frac{\rho\,\tau\,\Delta w}{\omega}.
\end{equation}
Here $E_{\text{BE}}$ is the minimum successful exit value required merely to break even in expected present value on the salary cut alone. Notice the two corollaries of~\eqref{eq:salary-equity-tradeoff}:
\begin{itemize}
  \item Taking a full market executive salary ($w_{\text{act}} = w_{\text{mkt}} \implies \Delta w = 0$) contributes zero sweat-capital subsidy, which is why a hired VP of Engineering paid \$240{,}000 in cash receives a $1\%$ to $2\%$ incentive option grant rather than a $30\%$ co-founder stake.
  \item Conversely, taking a steep salary cut ($\Delta w \gg 0$) is economically rational \emph{only} when paired with large founder-scale ownership $\omega$ and uncapped upside governance. Taking founder-level low pay for hired-employee single-digit or low-teen minority equity violates~\eqref{eq:salary-equity-tradeoff} on Day~1.
\end{itemize}

\section{Contracts When Partnering With an Existing Business}
\label{sec:existing-business-contracts}

Now apply everything in this chapter to a high-stakes contract arena many senior engineers encounter: an \textbf{established profitable business} (or an operating group controlling distribution, merchant relationships, regulatory licenses, or capital -- call it \textbf{ParentCo}) invites you to become ``Co-Founder / CEO / CTO / Managing Partner'' of a new venture or spin-out subsidiary (\textbf{NewCo}) with an equity stake of $10\%$ to $25\%$.

Unlike a standalone VC-backed startup -- where venture funds are passive minority preferred holders whose only path to liquidity is an eventual IPO or sale alongside common -- an operating ParentCo typically owns $75\%$ to $90\%$ of NewCo, controls NewCo's board of directors, transacts commercially with NewCo every month, and may prefer holding NewCo forever as a cash-generating division rather than ever selling it. Auditing such an offer requires checking four legal boundaries and running the exact expected-value math of a subsidiary case study.

\subsection{Which Entity and Which Instrument Are You Receiving?}
\label{subsec:entity-and-instrument}

First verify \textbf{which legal entity holds the core assets}. Confirm that NewCo itself -- the exact corporation or LLC issuing your equity -- owns the source code, customer contracts, trademarks, and payment or operating licenses (or holds an irrevocable, royalty-free perpetual license from ParentCo). Owning $20\%$ of a staffing shell whose software and customer contracts sit inside ParentCo gives you $20\%$ of an empty box.

Second, identify the exact legal instrument behind the words ``$20\%$ equity'':
\begin{enumerate}
  \item \textbf{Issued Voting Common Shares in NewCo}: true property ownership carrying statutory stockholder inspection rights, corporate fiduciary duties against majority self-dealing, voting power, and capital-gains treatment upon sale.
  \item \textbf{Non-Voting Shares}: pro-rata economic rights to declared dividends and sale proceeds, zero governance vote (viable only when paired with explicit minority veto covenants in a binding Shareholders' Agreement).
  \item \textbf{LLC Profits Interest Units}: when NewCo is a pass-through LLC already holding contributed assets worth $H > 0$ (say $H = \text{USD } 2{,}000{,}000$), granting an outright $20\%$ capital interest triggers immediate ordinary tax on \$400{,}000. Instead, a \textbf{Profits Interest} sets a \textbf{distribution hurdle} $H = \text{USD } 2{,}000{,}000$ so ParentCo receives the first \$2{,}000{,}000 upon liquidation and your units receive $20\%$ of operating profits and exit value above $H$ (yielding \$0 grant-date tax under IRS Rev.\ Proc.\ 93-27 / 2001-43, backed by a protective 30-day Section~83(b) filing). Because LLCs pass taxable income through on Schedule~K-1 even when zero cash is paid out, always require a \textbf{mandatory quarterly tax distribution clause} paying enough cash to cover K-1 taxes at the highest marginal rate.
  \item \textbf{Phantom Equity, Stock Appreciation Rights (SARs), or Contractual Revenue/Profit Share}: an unsecured corporate bonus promise taxed as ordinary income plus payroll tax, carrying zero shareholder inspection rights and vulnerable if ParentCo migrates customers or Version~2 products into a sister entity outside the bonus definition.
\end{enumerate}

\subsection{The Related-Party Cash-Extraction and Transfer-Pricing Trap}
\label{subsec:related-party-cash-extraction}

Suppose you hold genuine Voting Common Shares representing $20.00\%$ of NewCo, while ParentCo owns $80.00\%$ and controls the board. By Year~3, NewCo earns \$5{,}000{,}000 in external revenue against \$2{,}000{,}000 in third-party cloud, processing, and engineering costs -- a true standalone operating profit (EBITDA) of \textbf{\$3{,}000{,}000 per year}. Table~\ref{tab:related-party-extraction} shows how an unguarded parent--subsidiary structure sweeps $100\%$ of that profit out of your reach.

\begin{table}[htbp]
\centering
\caption{How unguarded related-party transfer pricing sweeps $100\%$ of NewCo's profit to ParentCo.}
\label{tab:related-party-extraction}
\footnotesize
\setlength{\tabcolsep}{5pt}
\begin{tabular}{@{}lrr@{}}
\toprule
\textbf{NewCo Annual Income Statement Line} & \textbf{Arm's-Length NewCo} & \textbf{ParentCo Sweep} \\
\midrule
External Customer Revenue & \$5{,}000{,}000 & \$5{,}000{,}000 \\
Third-Party Cloud, Processing \& Payroll Costs & (\$2{,}000{,}000) & (\$2{,}000{,}000) \\
\midrule
\textbf{Standalone Operating Profit (Before Parent Fees)} & \textbf{\$3{,}000{,}000} & \textbf{\$3{,}000{,}000} \\
ParentCo ``Shared Management \& Executive Fee'' & \$0 & (\$900{,}000) \\
ParentCo ``Brand, Regulatory \& IP License Royalty'' ($24\%$) & \$0 & (\$1{,}200{,}000) \\
ParentCo Marked-Up FX / Supply Clearing \& Director Fees & \$0 & (\$900{,}000) \\
\midrule
\textbf{NewCo Reported Pre-Tax Net Income} & \textbf{\$3{,}000{,}000} & \textbf{\$0} \\
\textbf{Cash Captured Directly by ParentCo ($0\%$ CTO share)} & \$0 & \textbf{\$3{,}000{,}000} \\
\textbf{Distributable Dividend Pool ($80 / 20$ split)} & \$3{,}000{,}000 & \textbf{\$0} \\
\textbf{Engineer CTO's $20\%$ Economic Share} & \textbf{\$600{,}000 / yr} & \textbf{\$0 / yr} \\
\bottomrule
\end{tabular}
\end{table}

\begin{trapbox}
\textbf{Related-Party Profit Sweeps in Majority-Owned Subsidiaries.}
In the right-hand column of Table~\ref{tab:related-party-extraction}, ParentCo's board majority causes NewCo to sign intercompany management, brand-royalty, and marked-up supply agreements with ParentCo ($100\%$ owned by the Majority Group, $0\%$ owned by you). Every dollar of NewCo's \$3{,}000{,}000 operating profit is booked as a deductible ``expense'' paid to ParentCo. NewCo reports \$0 Net Income, declares \$0 in dividends, and never needs to be sold because ParentCo already collects \$3{,}000{,}000 in cash annually above the Net Income line.

Closing this leak requires three binding Shareholders' Agreement covenants (Figure~\ref{fig:parentco-newco-boundary}):
\begin{enumerate}
  \item \textbf{Arm's-Length Related-Party Covenant and Minority Veto}: NewCo may not enter into, amend, or pay any fee, royalty, supply markup, loan, or executive charge to ParentCo or any affiliate (beyond pre-scheduled shared services billed strictly at documented third-party pass-through cost with zero markup) without \textbf{your prior written consent}.
  \item \textbf{Audit and General-Ledger Inspection Right}: quarterly financials, annual CPA-reviewed statements, and full access to intercompany contracts and invoices.
  \item \textbf{Defined Free-Cash-Flow Dividend Policy}: once NewCo holds a defined working-capital reserve (e.g., six months of operating expenses), at least $40\%$ to $50\%$ of excess annual Free Cash Flow must be distributed pro rata to all shareholders.
\end{enumerate}
\end{trapbox}

\begin{figure}[htbp]
\centering
\begin{tikzpicture}[x=1cm, y=1cm, >=Latex, font=\small]
  % External Customers box (left)
  \draw[rounded corners=4pt, fill=sourceblue!10, draw=sourceblue!80, thick] (0.0, 1.45) rectangle (2.55, 3.05);
  \node[align=center, font=\footnotesize\bfseries, text=sourceblue!90!black] at (1.275, 2.25) {External\\Customers\\[-1pt]\normalfont\scriptsize \$5.0M Revenue};

  % NewCo boundary box (center)
  \draw[rounded corners=5pt, fill=targetgreen!7, draw=targetgreen!75!black, very thick] (3.55, 0.70) rectangle (8.35, 3.80);
  \node[font=\footnotesize\bfseries, text=targetgreen!65!black] at (5.95, 3.48) {NewCo Operating Subsidiary};
  \node[font=\scriptsize, align=center] at (5.95, 3.05) {80\% ParentCo $\mid$ 20\% Engineer CTO};
  \draw[rounded corners=3pt, fill=white, draw=targetgreen!60] (3.85, 1.95) rectangle (8.05, 2.75);
  \node[font=\scriptsize, align=center] at (5.95, 2.35) {\textbf{Must Reside Inside NewCo:}\\Code IP $\cdot$ Customer Contracts $\cdot$ Licenses};
  \node[font=\scriptsize, align=center, text=black!80] at (5.95, 1.30) {Standalone Profit: \$3.0M\\[-1pt]\textbf{Guard 1:} Minority Veto on Affiliate Deals\\[-1pt]\textbf{Guard 2:} Mandatory FCF \& Tax Dividend Rule};

  % ParentCo box (top right)
  \draw[rounded corners=4pt, fill=trapred!8, draw=trapred!80, thick] (9.45, 2.55) rectangle (12.40, 4.05);
  \node[align=center, font=\footnotesize\bfseries, text=trapred!85!black] at (10.925, 3.30) {ParentCo\\[-1pt]\normalfont\scriptsize (100\% Majority Group,\\[-2pt]\normalfont\scriptsize 0\% Engineer CTO)};

  % Engineer CTO box (bottom right)
  \draw[rounded corners=4pt, fill=structurepurple!10, draw=structurepurple!80, thick] (9.45, 0.45) rectangle (12.40, 1.85);
  \node[align=center, font=\footnotesize\bfseries, text=structurepurple!85!black] at (10.925, 1.15) {Engineer CTO\\[-1pt]\normalfont\scriptsize 20\% Pro-Rata\\[-2pt]\normalfont\scriptsize Dividends \& Exit};

  % Arrows
  \draw[->, very thick, sourceblue!85] (2.55, 2.25) -- (3.55, 2.25);

  % Blocked / Guarded Related-Party Leak arrow (NewCo -> ParentCo)
  \draw[->, very thick, dashed, trapred] (8.35, 3.35) -- (9.45, 3.35);
  \node[font=\tiny\bfseries, text=trapred, align=center, fill=white, inner sep=1pt] at (8.90, 3.82) {Related-Party\\Fees / Markups};
  \node[font=\scriptsize\bfseries, text=trapred] at (8.90, 3.35) {$\times$};

  % Legitimate 80% dividend / exit flow (NewCo -> ParentCo)
  \draw[->, thick, targetgreen!75!black] (8.35, 2.75) -- (9.45, 2.85)
    node[midway, below, sloped, font=\tiny\bfseries] {80\% Pro-Rata};

  % Legitimate 20% dividend / exit flow (NewCo -> CTO)
  \draw[->, thick, targetgreen!75!black] (8.35, 1.40) -- (9.45, 1.25)
    node[midway, above, sloped, font=\tiny\bfseries] {20\% Pro-Rata};
\end{tikzpicture}
\caption{Entity and cash-flow boundary between ParentCo and NewCo. Without a minority veto over related-party charges, ParentCo can sweep $100\%$ of NewCo's operating profit via management fees and transfer-pricing markups (dashed red path), bypassing the $80 / 20$ shareholder waterfall.}
\label{fig:parentco-newco-boundary}
\end{figure}

\subsection{Good Leaver vs.\ Bad Leaver Repurchase Clauses}
\label{subsec:good-leaver-bad-leaver}

Private-subsidiary Shareholders' Agreements often include \textbf{Leaver Call Options} allowing ParentCo to repurchase your \emph{already-vested} shares when employment ends. Watch for drafts that define \emph{Good Leaver} solely as death, disability, or retirement, classifying voluntary resignation after Year~3 or \emph{termination without Cause} by ParentCo's board as a \emph{Bad Leaver} triggering repurchase of all vested shares at nominal par value (\$1 or \$40 total). Under such a clause, ParentCo can fire you without Cause the week before a \$30M sale and confiscate your vested shares for \$40.

Insist on the clean market rule: \emph{unvested} shares revert at original cost; \emph{vested} shares are either kept outright or repurchased strictly at \textbf{undiscounted independent Fair Market Value} upon resignation or termination without Cause, with cost-basis forfeiture restricted exclusively to narrow, objective \textbf{Cause} (proven fraud, embezzlement, felony conviction harming NewCo, or willful material breach uncured 30 days after written notice).

\subsection{Pre-Emptive Rights, Wash-Out Dilution, Tag-Along, and Drag-Along}
\label{subsec:preemptive-tag-drag}

Three further clauses protect minority shares from structural squeeze-outs:
\begin{enumerate}
  \item \textbf{Pre-Emptive (Pro-Rata) Rights and Insider FMV Pricing}: if NewCo has $1{,}000{,}000$ shares ($800{,}000$ ParentCo, $200{,}000$ CTO $= 20\%$) and is worth \$10{,}000{,}000 (\$10.00/share), an unguarded board could issue $19{,}000{,}000$ new shares exclusively to ParentCo at \$0.01 per share (\$190{,}000 internal cash), diluting your $20.00\%$ stake to $200{,}000 / 20{,}000{,}000 = 1.00\%$. Require both a \textbf{Pre-Emptive Right} to buy your $20\%$ pro-rata share of any new issuance on identical terms \emph{and} a rule that insider/affiliate issuances without minority consent must occur at independent Fair Market Value.
  \item \textbf{Tag-Along (Co-Sale) Right}: if ParentCo sells its $80\%$ block to an outside buyer, you have the right to sell your $20\%$ alongside ParentCo at the identical per-share price and terms rather than being stranded under a stranger.
  \item \textbf{Drag-Along Guardrails}: if ParentCo compels you to join a $100\%$ sale, require that the buyer be an \emph{unaffiliated third party}, consideration per share be identical, indemnity liability be capped at your net cash proceeds received, and no new post-sale non-compete be imposed without your separate consent.
\end{enumerate}

\subsection{Case Study: Joining an Existing Business's New Subsidiary as Founder or CTO}
\label{subsec:case-study-subsidiary-founder}

Let us combine the salary--equity model of Section~\ref{sec:founder-salary-vs-equity} with the subsidiary governance rules above on a concrete decision. Suppose your open-market engineering compensation is $w_{\text{mkt}} = \text{USD } 220{,}000$ per year ($4 \times 220{,}000 = \text{USD } 880{,}000$ over a $\tau = 4$-year build). Assume any new venture in this vertical has probability $p_s = 0.35$ of reaching a clean standalone equity value of $E_{\text{com}} = \text{USD } 20{,}000{,}000$ at Year~$T = 4$ (and probability $0.65$ of stalling at zero common payout), with opportunity discount rate $r = 10\%$ per year, so $(1 + r)^4 = 1.4641$ and equation~\eqref{eq:risk-multiplier} gives risk-and-illiquidity multiplier $\rho = 1.4641 / 0.35 = 4.1831$.

Compare three real-world offers across Table~\ref{tab:subsidiary-case-study}:
\begin{itemize}
  \item \textbf{Structure A (Independent VC-Backed Co-Founder)}: you and a co-founder start an independent C-Corp ($50\%$ each initially), taking \$40k in bootstrapped Year~1 and \$150k/year across Years~2--4 post-Seed/Series~A (4-year cash of \$490{,}000; cumulative forgone salary $\tau\Delta\bar{w} = \text{USD } 390{,}000$), diluting to $\omega = 35.0\%$ common ownership by Year~4 with board co-control.
  \item \textbf{Structure B (Balanced Subsidiary Partner / CTO)}: ParentCo owns $85\%$ of NewCo, supplies initial working capital and distribution, pays you a near-market base salary $w_{\text{act}} = \text{USD } 195{,}000$/year (4-year cash of \$780{,}000; modest forgone salary $\tau\Delta w = \text{USD } 100{,}000$), and grants $\omega = 15.0\%$ voting/profits equity protected by an arm's-length intercompany covenant, a $45\%$ excess-FCF dividend waterfall, and a Year-5 independent FMV buyout put option.
  \item \textbf{Structure C (Asymmetric Unpaid / Low-Pay Subsidiary Trap)}: ParentCo retains $90\%$ ownership and $100\%$ board control, pitches ``take a founder's sacrifice until NewCo is profitable'' at $w_{\text{act}} = \text{USD } 50{,}000$/year (4-year cash of \$200{,}000; forgone cash $\tau\Delta w = \text{USD } 680{,}000$) for $\omega = 10.0\%$ subsidiary equity with no intercompany fee cap and no buyout put right.
\end{itemize}

\begin{table}[htbp]
\centering
\caption{Four-year expected-value comparison ($w_{\text{mkt}}=\text{USD }220\text{k/yr}$, $p_s=0.35$, $r=10\%$, $\rho=4.1831$).}
\label{tab:subsidiary-case-study}
\footnotesize
\setlength{\tabcolsep}{2.8pt}
\begin{tabular}{@{}lrrr@{}}
\toprule
\textbf{Economic Metric (4-Year Horizon)} & \textbf{A: Independent} & \textbf{B: Balanced Sub.} & \textbf{C: Low-Pay Sub.} \\
\midrule
4-Year Cash Salary Received & \$490{,}000 & \$780{,}000 & \$200{,}000 \\
Cumulative Forgone Salary ($\tau\,\Delta w$) & \$390{,}000 & \$100{,}000 & \$680{,}000 \\
Year-4 Equity Ownership ($\omega$) & $35.0\%$ & $15.0\%$ & $10.0\%$ \\
Payout in \$20M Exit ($\omega \times \text{USD } 20\text{M}$, clean) & \$7{,}000{,}000 & \$3{,}000{,}000 & \$2{,}000{,}000 (or \$0)$^\dagger$ \\
Discounted Expected Equity PV ($\omega \times \text{USD } 20\text{M} / \rho$) & \$1{,}673{,}383 & \$717{,}164 & \$478{,}109 (or \$0)$^\dagger$ \\
\textbf{Net Expected Surplus vs.\ Market Job} & \textbf{+\$1{,}283{,}383} & \textbf{+\$617{,}164} & \textbf{--\$201{,}891 (to --\$680k)} \\
\textbf{Break-Even Successful Exit ($E_{\text{BE}}$)} & \textbf{\$4{,}658{,}000} & \textbf{\$2{,}789{,}000} & \textbf{\$28{,}445{,}000} \\
\bottomrule
\end{tabular}

\vspace{2pt}
{\scriptsize $^\dagger$If ParentCo sweeps profit via management fees and never sells NewCo, Structure~C's realized equity payout is \$0.}
\end{table}

Look at what Table~\ref{tab:subsidiary-case-study} proves quantitatively. Both Structure~A (independent founder: steep early pay cut paired with $35\%$ ownership and board voice) and Structure~B (subsidiary CTO: $15\%$ protected minority equity paired with near-market \$195k cash pay and a Year-5 FMV put exit) generate substantial positive expected surplus, breaking even at modest exit values below \$5M.

Structure~C fails basic expected-value arithmetic even \emph{before} any governance abuse: sacrificing \$680{,}000 of guaranteed cash for $10\%$ of a risky subsidiary requires a \$28{,}445{,}000 clean exit merely to tie your day job, so even a \$20{,}000{,}000 win leaves you \$201{,}891 poorer in expected present value -- effectively donating unpaid engineering R\&D to subsidize ParentCo's $90\%$ majority block. Whenever an existing business retains majority control ($70\%$ to $90\%$) of a new subsidiary, insist on four aligned terms: (1) near-market cash base pay (or an automatic contractual step-up to full market salary once defined funding or revenue milestones hit), (2) arm's-length intercompany pricing and management-fee caps, (3) a mandatory pro-rata profit/dividend distribution waterfall while held privately, and (4) an independent Fair Market Value buyout put/call window after Year~4 or Year~5 if no external sale has occurred.

\subsection{Ten-Point Contract Diagnostic Checklist}
\label{subsec:ten-point-contract-diagnostic}

Before signing any co-founder, early-executive, or corporate-subsidiary equity agreement, verify every line of Table~\ref{tab:ten-point-diagnostic} against the binding legal documents.

\begin{table}[htbp]
\centering
\caption{Ten-point diagnostic checklist for founder and existing-business equity contracts.}
\label{tab:ten-point-diagnostic}
\footnotesize
\setlength{\tabcolsep}{4.5pt}
\begin{tabular}{@{}cp{2.55cm}p{4.35cm}p{4.55cm}@{}}
\toprule
\textbf{\#} & \textbf{Checkpoint} & \textbf{Contract Red Flag} & \textbf{Clean Protective Standard} \\
\midrule
1 & Entity \& IP Boundary & ParentCo owns IP, brand, or key licenses on revocable terms. & Issuing entity (NewCo) owns code, contracts, and licenses outright. \\
\addlinespace[2pt]
2 & Instrument Class & Unsecured ``phantom bonus'' or undefined ``$20\%$ equity.'' & Voting Common Stock or hurdle LLC Profits Interest + cap table exhibit. \\
\addlinespace[2pt]
3 & Salary--Equity Balance & Unpaid/steep pay cut for $10\%$--$15\%$ parent-controlled subsidiary stake. & Founder-scale equity ($35\%+$) for sweat cuts, or near-market pay in a subsidiary. \\
\addlinespace[2pt]
4 & Mutual Vesting & One founder unvested while the other holds $100\%$ Day-1 stock. & Symmetric 4-yr reverse vesting, 1-yr cliff, credit for prior full-time months. \\
\addlinespace[2pt]
5 & Tax Elections \& K-1s & Missed 30-day 83(b) window; LLC allocates K-1 tax with \$0 cash. & File 83(b) within 30 calendar days; mandatory LLC quarterly tax distributions. \\
\addlinespace[2pt]
6 & Sale Acceleration & Zero acceleration on post-merger firing, or $100\%$ single-trigger. & Double-trigger acceleration (sale + termination without Cause / Good Reason). \\
\addlinespace[2pt]
7 & Liquidation Stack & $>1\times$ senior preference, uncapped double dip, or full ratchet. & $1\times$ non-participating preferred; broad-based weighted-average anti-dilution. \\
\addlinespace[2pt]
8 & Related-Party Leak & Majority owner can bill management fees, royalties, or markups. & Written minority veto over any ParentCo / affiliate contract or fee above cost. \\
\addlinespace[2pt]
9 & Leaver \& Exit Put & Resignation or no-cause firing forfeits \emph{vested} shares; zero exit path. & Vested shares kept or bought at undiscounted FMV; Year-5 FMV put in subsidiaries. \\
\addlinespace[2pt]
10 & Dilution, Tag \& Drag & Insider cheap-share washout; drag-along to affiliate or uncapped liability. & Pro-rata pre-emptive right + FMV insider rule; Tag-Along; 3rd-party Drag cap. \\
\bottomrule
\end{tabular}
\end{table}

With corporate ownership ledgers, preferred waterfalls, and minority shareholder contracts established, Part~IV turns to the operational plumbing of money movement itself: mobile-money unit economics, customer float accounting, foreign-exchange liquidity corridors, stablecoin settlement rails, and bank reconciliation controls.

\section*{Exercises}
\addcontentsline{toc}{section}{Exercises}

\begin{exercise}[Non-Participating vs.\ Participating Exit Waterfall]
A startup raises a single venture round of $I = \text{USD } 8{,}000{,}000$ for $\alpha = 25.0\%$ post-money ownership ($V_{\text{post}} = \text{USD } 32{,}000{,}000$), with Common Stock holding the remaining $75.0\%$. Compute the exact dollar payouts $(P_{\text{inv}}, P_{\text{com}})$ at three acquisition exit values -- \$6M, \$20M, and \$48M -- under: (a) $1\times$ Non-Participating Preferred; (b) $1\times$ Uncapped Participating Preferred; and (c) $1\times$ Participating Preferred capped at $2.5\times$ invested capital (\$20{,}000{,}000 cap). State the exact conversion threshold $E^*$ for parts~(a) and~(c).
\end{exercise}

\begin{exercise}[Stacked Seniority vs.\ Pari Passu]
Series~A invested \$6{,}000{,}000 ($1\times$ non-participating, $\alpha_A = 20\%$) and Series~B invested \$18{,}000{,}000 ($1.5\times$ non-participating, so $L_B = \text{USD } 27{,}000{,}000$, $\alpha_B = 25\%$), with Common holding $55\%$. Suppose the company sells for $E = \text{USD } 22{,}000{,}000$. Compute the dollar payout to Series~B, Series~A, and Common if: (a) Series~B is strictly senior to Series~A; and (b) Series~B and Series~A rank \emph{pari passu} in proportion to their preference amounts $(L_B, L_A)$. At what minimum exit value $E$ does Common receive its first positive dollar?
\end{exercise}

\begin{exercise}[Broad-Based Weighted Average vs.\ Full Ratchet]
Prior to a down round, a startup has $6{,}000{,}000$ common shares, $1{,}000{,}000$ reserved options, and $3{,}000{,}000$ Series~A Preferred shares originally purchased at $CP_1 = \text{USD } 2.50$ per share (\$7{,}500{,}000 total Series~A investment), giving $A = 10{,}000{,}000$. The company raises $I_{\text{down}} = \text{USD } 2{,}000{,}000$ in Series~B at $p_{\text{down}} = \text{USD } 1.25$ per share. Compute the adjusted Series~A conversion price $CP_2$ (to four decimal places), Series~A's new as-converted common share count, and the founders' post-Series-B fully diluted ownership percentage under (a) Broad-Based Weighted Average anti-dilution and (b) Full Ratchet anti-dilution.
\end{exercise}

\begin{exercise}[The Cost of Missing the 30-Day Section 83(b) Deadline]
A founder buys $3{,}600{,}000$ common shares on Day~1 at \$0.00001 par value (\$36 cash) on a 4-year reverse-vesting schedule with a 1-year cliff, and misses the 30-day Section~83(b) deadline. At Month~12 (the cliff), the company's 409A Common FMV is \$0.30 per share. Across Months~13--24, 409A FMV is \$0.60 per share. At the end of Month~24, the founder departs amicably (keeping vested shares while the company repurchases unvested shares at par). Assuming a $38\%$ combined marginal ordinary-income tax rate, compute the cash tax owed across Years~1 and~2 before selling a single share, and compare it with the tax owed at Months~12 and~24 had the 83(b) election been filed on Day~5.
\end{exercise}

\begin{exercise}[Salary--Equity Break-Even in a Corporate Spin-Out]
An engineer earning $w_{\text{mkt}} = \text{USD } 240{,}000$ per year is offered the CEO role of a new fintech subsidiary $82\%$ owned by an established bank ($18\%$ voting common stock to the CEO, protected by arm's-length covenants). Assuming a $\tau = 4$-year build, exit/put horizon $T = 4$ years, annual discount rate $r = 10\%$, and probability $p_s = 0.40$ of achieving a successful standalone valuation, compute the risk-and-illiquidity multiplier $\rho$ and the minimum successful subsidiary valuation $E_{\text{BE}}$ required for the CEO to break even in expected present value if the subsidiary pays annual cash salary: (a) $w_{\text{act}} = \text{USD } 200{,}000$; (b) $w_{\text{act}} = \text{USD } 80{,}000$; and (c) $w_{\text{act}} = 0$.
\end{exercise}

\begin{exercise}[Diagnosing a Spin-Out Term Sheet With an Existing Business]
ParentCo (a profitable logistics company) invites you to build an automated freight-matching platform inside a new entity, FreightTech LLC, offering: (i) $12\%$ ``Phantom Profit Units'' paying $12\%$ of FreightTech LLC's GAAP Net Income upon an eventual sale of FreightTech LLC, paired with a \$60{,}000 annual base salary (versus your \$210{,}000 market wage); (ii) ParentCo retains ownership of the freight software repository and licenses it to FreightTech LLC for an annual fee set each January by FreightTech LLC's board (controlled $3\text{--}0$ by ParentCo); and (iii) if your employment terminates for any reason prior to a sale, all Phantom Profit Units expire for \$0. Identify five distinct structural defects in this proposal using Sections~\ref{sec:founder-salary-vs-equity}--\ref{sec:existing-business-contracts}, and rewrite (i), (ii), and (iii) into clean, bankable counterparty terms.
\end{exercise}
